\documentclass[11pt]{article}

\usepackage[margin=1in]{geometry}
\usepackage{booktabs}
\usepackage{longtable}
\usepackage{array}
\usepackage{enumitem}
\usepackage{hyperref}
\hypersetup{colorlinks=true, linkcolor=blue, urlcolor=blue}

\title{Continuation Value is All You Need: \\
Computational Methods for Heterogeneous-Agent Macro}
\author{Jeffrey Sun$^1$\footnote{University of Toronto. Based on a course originally taught at Zhejiang University made possible by a CES--Chow Fellowship for Short-Term Teaching in China.}}
\date{}

\begin{document}
\date{May 2026}
\maketitle

\vspace*{-1em}
\section*{Course Overview}

We study computational methods for solving macroeconomic models with dynamics, heterogeneous agents, and varying
assumptions about expectations. The course is interactive with Jupyter notebooks written in the programming language Julia.

The code introduces my package HouseholdStages.jl for building very rich heterogeneous-agent household blocks.

Starting with a simple consumption-savings
model of a single agent, we will build up to representative-agent
general equilibrium models, then heterogeneous-agent general
equilibrium models. We then show how to solve these models
in steady state, in deterministic transitions, and two methods for
solving them with aggregate uncertainty: the Krusell-Smith method
and my own Neural Value Function Iteration (nVFI).

The course consists of 8 lectures, each of which has a
slide deck and Julia notebook. The notebooks are not quite standalone.
For self-study, the ideal reading order is to read the slides first
before running the corresponding notebook.

\section*{Lectures}

\begin{enumerate}
  \item \textbf{Introduction.} The consumption--savings problem in recursive form. The Bellman equation.

  \item \textbf{Computation.} The Julia language. Implementing Value Function Iteration (VFI).

  \item \textbf{Heterogeneity.} Simulating individuals and populations forward given the value function.

  \item \textbf{Equilibrium.} Heterogeneous agents in equilibrium: the Aiyagari model.  The Household Stage concept and HouseholdStages.jl package.

  \item \textbf{Expectations.} Transition dynamics in the Aiyagari model.

  \item \textbf{Spatial Heterogeneity and Calibration.} A spatial model with Gumbel migration. An example of indirect inference: calibrating location preference shifters to match populations.

  \item \textbf{Uncertainty.} The Krusell--Smith model. The Krusell--Smith method.

  \item \textbf{Continuation Value is All You Need.} Solving heterogeneous-agent model with aggregate uncertainty using Neural Value Function Iteration.
\end{enumerate}

\end{document}
